\section{方法与实验协议}
\label{sec:method}

\subsection{预算受限主动重观测}
\label{sec:baro}
语言导航闭环给出候选建筑后，决策问题收窄为：当前观测的校准概率是否足以定类，若否，是否花一次额外航时去下降并居中。
动作集为
\begin{equation}
\calA=\{\mathrm{hold},\,\mathrm{descend\_center}\}.
\end{equation}
$\mathrm{hold}$ 保持巡航高度继续使用当前裁块；$\mathrm{descend\_center}$ 先平移到目标估计位置正上方，再下降固定高度步长（默认 $10\,\mathrm{m}$），在更细的有效 GSD 上重跑损伤头。
平均每样本允许 $B$ 次额外原生档观测；每位置与每回合次数有硬上限（同一位置最多 2 次复检）。

给定四类损伤的校准概率 $p_t\in\Delta^{K-1}$，$K=4$，定义归一化熵
\begin{equation}
\label{eq:entropy}
U_t=-\frac{1}{\log K}\sum_{k=1}^{K}p_{t,k}\log p_{t,k}\in[0,1].
\end{equation}
下降动作的预期收益不是手调高度比，而是在验证裁块上拟合的条件期望熵减
\begin{equation}
\label{eq:ig}
\widehat{\mathrm{IG}}(a)=U_t-\widehat{\mathbb{E}}[U\mid \mathrm{GSD}(a),\hat y_t],
\end{equation}
其中 $\mathrm{GSD}(a)$ 由高度--分辨率阶梯给出，$\hat y_t=\arg\max_k p_{t,k}$。
$\widehat{\mathrm{IG}}$ 是查找表意义上的期望熵减，不是 Shannon 信息增益的无偏估计，不得写成贝叶斯最优。
并列方法用分拆共形预测：当预测集合 $|C(x)|>1$ 时认为证据不足以唯一定类，从而触发重观测，并报告 $1-\alpha$ 覆盖。

\subsection{高度--有效 GSD 分辨率阶梯}
\label{sec:gsd}
xBD 只有单一原生分辨率（约 $0.5\,\mathrm{m/px}$）。
若下降只是裁小窗口再把分类裁块双线性缩放到固定 $96\times96$，则不引入新信息。
为此定义线性阶梯：巡航高度 $30\,\mathrm{m}$ 对应 $4\times$ 降采样加高斯模糊，最低高度 $10\,\mathrm{m}$ 对应原生像素，中间高度内插尺度。
有效 GSD 写入每步观测日志，使式~\eqref{eq:ig} 的条件期望成为可拟合的量。

该构造的适用范围必须写明：仿真中的``下降''等于逐步撤销我们自己加上的模糊，最好情况下测到的是分类器能否从合成降质中恢复，而不是真实 $10\,\mathrm{m}$ 飞行所提供的厘米级新纹理。
因此后文零结果只约束``在该有效 GSD 阶梯上机制是否可识别''，不能外推为``真实降高无信息''。

\subsection{双时相损伤分类与校准}
\label{sec:classifier}
共享编码器分别提取灾前/灾后特征，融合后接四类损伤头与二值变化头（多任务）。
两种融合方式并列评测（预注册 H2）：\textbf{拼接}（$[f_{\mathrm{pre}},f_{\mathrm{post}},f_{\mathrm{post}}-f_{\mathrm{pre}}]$）与
\textbf{差分注意力}（D2ANet 式门控：先对 $[f_{\mathrm{pre}},f_{\mathrm{post}}]$ 学通道门控再相减，再对差分特征学第二层门控）。
正文实验默认报告差分注意力档；\tabref{tab:cpv2} 给出两者的正式对比。
验证集上拟合标量温度 $T$，推理用 $\mathrm{softmax}(\mathrm{logits}/T)$。
温度缩放不改变 $\arg\max$，故期望校准误差（ECE）改善不得写成识别精度改善。
宏观 F1、分类型可靠性与总体 ECE 一并报告，以避免基率混淆。

\paragraph{类别加权与按宏观 F1 选择 checkpoint。}
训练集完好类占比 66\%--79\%，早期版本用未加权交叉熵训练、按验证集 overall accuracy 选最佳 epoch，
这一组合在该不平衡度下会系统性选出常数预测器。
修复为：(i) 交叉熵按训练集逐类频率倒数加权；(ii) 模型选择与早停均改用验证集宏观 F1；
(iii) 每次训练收尾对重新加载后的 checkpoint 现场核验宏观 F1 与 accuracy 与训练期间记录值一致，不一致则拒绝保存。
训练/验证曲线与逐类召回见 \secref{sec:x1train}。

\subsection{触发策略}
\label{sec:trigger}
\paragraph{条件期望熵。}
在验证裁块上按 GSD 档与预测类统计平均熵，得到查找表 $\widehat{\mathbb{E}}[U\mid\mathrm{GSD},\hat y]$。
在线触发当 $\widehat{\mathrm{IG}}(\mathrm{descend\_center})$ 超过阈值且预算未耗尽。
若表缺失则回退到高度比启发式，并在日志中标明。

\paragraph{共形集合。}
在验证集上用 Adaptive Prediction Sets 拟合分位数 $\hat q_{1-\alpha}$。
预测集合取累积概率达到 $\hat q$ 的最小标签子集；当 $|C(x)|>1$ 时触发下降。

对照包括：永不下降、随机花预算、未校准峰值/熵阈值、固定下降。
零触发回合对条件熵差视为缺失，触发率报 0，不得填补成功下降。

\subsection{事件不相交数据协议}
\label{sec:split}
灾害事件而非作物块是划分单位。训练、温度拟合（验证）、测试与留出使用互斥事件集合，见附录。
任一评测题的 \texttt{disaster} 落入训练或验证集合即视为泄漏并在生成脚本中断言失败。
题库审核字段记录为``模型辅助审核 + 作者抽查''，不得写成纯人工审核。

\subsection{运动学参数与预算换算}
\label{sec:kinematics}
语言导航默认步数预算 $12$，判到达半径 $35\,\mathrm{m}$，成功半径 $25\,\mathrm{m}$，
朝目标单步上限 $80\,\mathrm{m}$，未命中时探索步长 $90\,\mathrm{m}$。
复检每次下降 $10\,\mathrm{m}$，同一位置最多 $2$ 次，水平居中位移上限 $40\,\mathrm{m}$。
仿真器为实时演示把单段飞行时长封顶为 $8\,\mathrm{s}$，因此墙钟时间不能当作航时；
按距离/速度换算、忽略该封顶：探索步 $90\,\mathrm{m}$ 约 $7.5\,\mathrm{s}$，$12$ 步巡航约 $90\,\mathrm{s}$；
两次下降共 $20\,\mathrm{m}$ 约 $1.7\,\mathrm{s}$，一次 $40\,\mathrm{m}$ 居中约 $3.3\,\mathrm{s}$。
离线实验中平均每样本额外观测预算 $B=0.25$ 对应每四个裁块一次下降--居中，
相对 $12$ 步巡航预算约为数个百分点的等效航时，而不是电池或气动模型。
该换算只给出数量级，不能外推到真实续航。

\subsection{评价指标与可识别性判据}
\label{sec:metrics}
端到端到达指标如下，首次出现即给出定义，后文用缩写。
\textbf{严格成功率} \SR：回合终止时无人机与目标的水平距离 $\le 25\,\mathrm{m}$ 的比例。
\textbf{语义成功率} \semSR：在指令所述损伤类别的真值建筑集合上，距离最近真值 $\le 25\,\mathrm{m}$ 的比例；它宽于 \SR，不得写成严格成功。
\textbf{导航误差} \NE：终止位置到目标的水平距离（米），报告均值。
\textbf{路径效率} \SPL：成功回合按最短路径与实际路径长度比加权，失败回合记 0，再对题平均。

离线判断指标为总体精度、宏观 F1、ECE，以及预测翻转次数
\begin{equation}
n_{\mathrm{flip}}=\#\{i: \hat y_i^{\mathrm{cruise}}\neq \hat y_i^{\mathrm{native}}\}.
\end{equation}
可识别性判据预先写明，而不是为负结果事后找解释：
若 $n_{\mathrm{flip}}$ 过小，或校准策略与不下降、随机之间的宏观 F1 差距（相对逐样本翻转参照）含 0，
则任何预算分配策略在该数据上都\textbf{不可识别}——正确结论是机制未得到检验，而不是机制无效于导航。
逐样本翻转知识的贪心分配是宏观 F1 的参照而不是上界：宏观 F1 对样本不可加，策略曲线超过该参照并不构成逻辑矛盾，后文改称``oracle 参照''。
